% Generated draft: requires mathematical and source review.
\documentclass[11pt]{article}
\usepackage[margin=25mm]{geometry}
\usepackage{fontspec}
\setmainfont{Latin Modern Roman}
\usepackage{amsmath,amssymb,mathtools,unicode-math}
\setmathfont{Latin Modern Math}
\usepackage{enumitem,needspace,xurl,hyperref}
\setlength{\parindent}{0pt}
\setlength{\parskip}{5pt}
\setlength{\emergencystretch}{3em}
\newcommand{\N}{\mathbb{N}}
\newcommand{\Z}{\mathbb{Z}}
\newcommand{\Q}{\mathbb{Q}}
\newcommand{\R}{\mathbb{R}}
\newcommand{\C}{\mathbb{C}}
\newcommand{\F}{\mathbb{F}}
\newcommand{\A}{\mathbb{A}}
\newcommand{\PP}{\mathbb{P}}
\newcommand{\E}{\mathbb{E}}
\newcommand{\eps}{\varepsilon}
\newcommand{\dd}{\,\mathrm{d}}
\begin{document}
\section*{UnsolvedMath ID: 801043}
\textbf{Source problem code:} catalogue entry
\subsection{Definitions and mathematical statement}
Let \(F_\infty\) be the free group on countably infinitely many generators, and let \(\lambda\) be Lebesgue probability measure on \([0,1]\). Write
\[
X=\operatorname{Act}(F_\infty,([0,1],\lambda)).
\]
Under the usual probability-measure-preserving interpretation, an element \(x\in X\) is an action of \(F_\infty\) by invertible measurable transformations preserving \(\lambda\). The supplied record does not define the action space, its topology or standard Borel structure, or whether transformations are identified modulo equality almost everywhere.

For \(x\in X\), the source denotes its Schreier graph on \([0,1]\) by \(\operatorname{Sch}(x)\). Given a fixed free generating sequence \((s_n)_{n\in\N}\), a conventional undirected Schreier graph joins distinct vertices \(u,v\) when \(v=s_n\cdot_x u\) or \(u=s_n\cdot_x v\) for some \(n\). The source does not specify the generating sequence or its conventions concerning loops, multiple edges, and measurable representatives.

For a graph with vertex set \([0,1]\), a measurable perfect matching can be expressed as a measurable map \(m:[0,1]\to[0,1]\) such that, for every \(u\in[0,1]\),
\[
m(m(u))=u,\qquad m(u)\ne u,\qquad
\{u,m(u)\}\text{ is an edge}.
\]
A measure-theoretic variant requires these conditions only on a conull measurable set. Whether the source requires a matching everywhere or only almost everywhere, and whether measurability means Borel or Lebesgue measurability, remains unspecified.

With the intended conventions fixed, put
\[
\mathcal M=\{x\in X:\operatorname{Sch}(x)\text{ has a measurable perfect matching}\}.
\]
The question is whether \(\mathcal M\) is \(\boldsymbol{\Sigma}^{1}_{1}\)-complete. Under the standard Borel-reduction interpretation, this asks whether \(\mathcal M\) is analytic and, for every standard Borel space \(Y\) and every analytic set \(B\subseteq Y\), there exists a Borel map \(f:Y\to X\) satisfying
\[
\forall y\in Y,\qquad y\in B\ \Longleftrightarrow\ f(y)\in\mathcal M.
\]
The supplied record does not specify the reduction convention. Consequently, this is a provisional formulation: the action-space, graph, measurability, null-set, and reduction conventions must be confirmed before it determines a unique precise problem.
\subsection{Short English statement}
For actions of the countably generated free group on the Lebesgue probability interval, ask whether the actions whose Schreier graphs admit measurable perfect matchings form an analytic-complete set. The supplied statement leaves several conventions unresolved, especially matching everywhere versus almost everywhere and the meaning of completeness.
\subsection{Sources}
\begin{itemize}
\item Local UnsolvedMath registry: permanent ID \(801043\), exact source problem code \texttt{catalogue entry}, title \emph{Open Problems in Mathematical Logic --- Problem 55}. The supplied registry attributes the problem to Barmpalias, Gao et al., with proposed year 2026, and labels its status as open; that label is supplied metadata, not an independently verified mathematical assessment.
\item Barmpalias, Gao et al., \emph{Open Problems in Mathematical Logic} (2026), Problem 55. Supplied source URL: \url{https://arxiv.org/abs/2608.26628}. The record reports access on 2026-10-16 and extraction by \texttt{arxiv\_latex}. These are supplied provenance details; the source was not independently accessed or verified for this draft.
\end{itemize}
\end{document}
